
\documentclass{article}
\usepackage{colortbl}
\usepackage{makecell}
\usepackage{multirow}
\usepackage{supertabular}

\begin{document}

\newcounter{utterance}

\centering \large Interaction Transcript for game `cladder', experiment `full\_v1.5\_default', episode 5310 with BSU{-}SLIM\_\_Qwen3.5{-}9B{-}Base\_\_playpen{-}prm{-}9b{-}best\_of\_n.
\vspace{24pt}

{ \footnotesize  \setcounter{utterance}{1}
\setlength{\tabcolsep}{0pt}
\begin{supertabular}{c@{$\;$}|p{.15\linewidth}@{}p{.15\linewidth}p{.15\linewidth}p{.15\linewidth}p{.15\linewidth}p{.15\linewidth}}
    \# & \multicolumn{2}{c}{Player} && \multicolumn{2}{c}{Game Master} \\
    \hline

    \theutterance \stepcounter{utterance}  
    & & & \multicolumn{4}{p{0.6\linewidth}}{
        \cellcolor[rgb]{0.9,0.9,0.9}{
            \makecell[{{p{\linewidth}}}]{
                \texttt{\tiny{[P1$\langle$GM]}}
                \texttt{You are an expert in causal inference. The following question is not a typical commonsense query, but rather a meticulously designed question created by a professor specializing in causal inference, intended to assess the students' mastery of the course content.} \\
\texttt{Imagine a self{-}contained, hypothetical world with only the following conditions, and without any unmentioned factors or causal relationships: Liking spicy food has a direct effect on occupation and salary. Skill has a direct effect on occupation and salary. Occupation has a direct effect on salary. Skill is unobserved. Method 1: We look at how liking spicy food correlates with salary case by case according to occupation. Method 2: We look directly at how liking spicy food correlates with salary in general. To understand how liking spicy food affects salary, is it more correct to use the Method 1 than Method 2?} \\
\texttt{Start your answer with 'Yes' or 'No', followed by additional reasoning or evidence to support your explanation.} \\
            }
        }
    }
     \\ \\

    \theutterance \stepcounter{utterance}  
    & \multicolumn{4}{p{0.6\linewidth}}{
        \cellcolor[rgb]{0.561,0.737,0.561}{
            \makecell[{{p{\linewidth}}}]{
                \texttt{\tiny{[P1$\rangle$GM]}}
                \texttt{Yes} \\
\\ 
\texttt{**Reasoning:**} \\
\\ 
\texttt{In this hypothetical world, **Skill** acts as an unobserved confounder because it is a common cause of both **Spicy Food Liking** and **Salary**. To correctly estimate the *direct* (causal) effect of spicy food on salary, we must block the back{-}door paths that introduce bias. Specifically, we need to control for the confounding variable (Skill) to isolate the unique contribution of spicy food.} \\
\\ 
\texttt{*   **Analysis of Method 2 (Unadjusted Correlation):** This method looks at the raw correlation between spicy food and salary. Because skill influences both variables, the observed correlation will capture two distinct effects:} \\
\texttt{    1.  The effect of spicy food on salary.} \\
\texttt{    2.  The effect of skill on salary (which is correlated with spicy food).} \\
\texttt{    Consequently, Method 2 conflates the causal effect of spicy food with the effect of skill, leading to biased estimates of the causal effect of spicy food.} \\
\\ 
\texttt{*   **Analysis of Method 1 (Stratified Analysis by Occupation):** The prompt states that **Occupation** is a *direct effect* of spicy food (Spicy Food \$\textbackslash{}rightarrow\$ Occupation). In causal diagrams (DAGs), this means Occupation is a mediator. If we stratify (or condition) on Occupation, we block the path from Spicy Food to Salary that runs through Occupation.} \\
\\ 
\texttt{    However, the critical question is whether stratifying on Occupation also blocks the confounding path from Skill to Salary. Since Skill causes Occupation, and Occupation is a collider on the path from Spicy Food to Salary (via the arrow from Spicy Food), we must be careful.} \\
\\ 
\texttt{    Let's refine the causal structure based on the prompt:} \\
\texttt{    1.  Spicy Food \$\textbackslash{}rightarrow\$ Occupation} \\
\texttt{    2.  Skill \$\textbackslash{}rightarrow\$ Occupation} \\
\texttt{    3.  Skill \$\textbackslash{}rightarrow\$ Salary} \\
\texttt{    4.  Occupation \$\textbackslash{}rightarrow\$ Salary} \\
\\ 
\texttt{    The back{-}door paths from Spicy Food to Salary are:} \\
\texttt{    *   Path A: Spicy Food \$\textbackslash{}rightarrow\$ Salary (Direct causal path of interest).} \\
\texttt{    *   Path B: Spicy Food \$\textbackslash{}rightarrow\$ Occupation \$\textbackslash{}leftarrow\$ Skill \$\textbackslash{}rightarrow\$ Salary. (This is a back{-}door path because it enters the causal node Spicy Food through an arrow, goes through Occupation, and continues to Salary).} \\
\\ 
\texttt{    *   **Method 2** fails because it leaves Path B open, allowing Skill to bias the estimate.} \\
\\ 
\texttt{    *   **Method 1** conditions on Occupation. By conditioning on Occupation, we block Path B (Spicy Food \$\textbackslash{}rightarrow\$ Occupation \$\textbackslash{}leftarrow\$ Skill). Since Occupation is a *collider* (it receives arrows from both Spicy Food and Skill), conditioning on a collider typically opens a path between its causes. However, in this specific structure, we are interested in the effect of Spicy Food. By fixing the occupation, we remove the variation in occupation that is driven by skill. Since the only way Spicy Food influences Salary *other than* directly is through the channel of Skill (via Occupation), controlling for Occupation effectively blocks the confounding influence of Skill on the relationship between Spicy Food and Salary.} \\
\\ 
\texttt{    Therefore, Method 1 yields an unbiased estimate of the direct causal effect of spicy food on salary, whereas Method 2 remains biased due to the omitted variable bias caused by the unobserved Skill.} \\
            }
        }
    }
    & & \\ \\

    \theutterance \stepcounter{utterance}  
    & & & \multicolumn{2}{p{0.3\linewidth}}{
        \cellcolor[rgb]{0.9,0.9,0.9}{
            \makecell[{{p{\linewidth}}}]{
                \texttt{\tiny{[GM$|$GM]}}
                \texttt{Yes} \\
            }
        }
    }
    & & \\ \\

    \theutterance \stepcounter{utterance}  
    & & & \multicolumn{2}{p{0.3\linewidth}}{
        \cellcolor[rgb]{0.9,0.9,0.9}{
            \makecell[{{p{\linewidth}}}]{
                \texttt{\tiny{[GM$|$GM]}}
                \texttt{no} \\
            }
        }
    }
    & & \\ \\

    \theutterance \stepcounter{utterance}  
    & & & \multicolumn{2}{p{0.3\linewidth}}{
        \cellcolor[rgb]{0.9,0.9,0.9}{
            \makecell[{{p{\linewidth}}}]{
                \texttt{\tiny{[GM$|$GM]}}
                \texttt{game\_result = LOSE} \\
            }
        }
    }
    & & \\ \\

\end{supertabular}
}

\end{document}
